
\medskip
\phantomsection\label{def:sga3-VIB-11.8.0}%
\noindent\textbf{Definition 11.8.0.---}%
\footnote{\label{note:sga3-VIB-112}Editor's note: the
statements \SGARefLink{lem:sga3-VIB-11.8}{11.8} and \SGARefLink{prop:sga3-VIB-11.9}{11.9} concern only the notion of a
comodule over a coalgebra. We have therefore introduced
\hyperref[def:sga3-VIB-11.8.0]{Definition 11.8.0} and reformulated \hyperref[lem:sga3-VIB-11.8]{11.8} and \hyperref[prop:sga3-VIB-11.9]{11.9}
accordingly.}
Let \(S\) be a scheme. An \(\mathcal O_S\)-\emph{coalgebra}
is an \(\mathcal O_S\)-module \(\mathcal C\), together with
two morphisms of \(\mathcal O_S\)-modules
\[
  \Delta:\mathcal C\longrightarrow\mathcal C\otimes\mathcal C,
  \qquad
  \varepsilon:\mathcal C\longrightarrow\mathcal O_S,
\]
satisfying the following two axioms \textup{(cf.\ I, \SGARefLink{sga3:i:subsection:4-2}{4.2})}:
\begin{enumerate}
\item[\textup{(CO 1)}]\phantomsection\label{ax:sga3-VIB-11.8.0-CO1}
  \(\Delta\) is coassociative: the following diagram is
  commutative:
  \begin{center}
    \phantomsection\label{diag:sga3-VIB-11.8.0-coassociativity}%
    \begin{tikzcd}[row sep=large,column sep=large]
      & \mathcal C\otimes\mathcal C
        \arrow[dr,"\mathrm{id}\otimes\Delta"] & \\
      \mathcal C
        \arrow[ur,"\Delta"]
        \arrow[dr,"\Delta"'] & &
      \mathcal C\otimes\mathcal C\otimes\mathcal C \\
      & \mathcal C\otimes\mathcal C
        \arrow[ur,"\Delta\otimes\mathrm{id}"'] &
    \end{tikzcd}
  \end{center}

\item[\textup{(CO 2)}]\phantomsection\label{ax:sga3-VIB-11.8.0-CO2}
  \(\varepsilon\) is a counit; that is, the following two
  composites are the identity:
  \phantomsection\label{form:sga3-VIB-11.8.0-counit-right}%
  \[
    \mathcal C\xrightarrow{\Delta}
    \mathcal C\otimes\mathcal C
    \xrightarrow{\mathrm{id}\otimes\varepsilon}
    \mathcal C\otimes\mathcal O_S
    \xrightarrow{\sim}\mathcal C,
  \]
  \phantomsection\label{form:sga3-VIB-11.8.0-counit-left}%
  \[
    \mathcal C\xrightarrow{\Delta}
    \mathcal C\otimes\mathcal C
    \xrightarrow{\varepsilon\otimes\mathrm{id}}
    \mathcal O_S\otimes\mathcal C
    \xrightarrow{\sim}\mathcal C.
  \]
\end{enumerate}

A \emph{\(\mathcal C\)-comodule \textup{(on the right)}} is
an \(\mathcal O_S\)-module \(\mathcal E\), together with a
morphism of \(\mathcal O_S\)-modules
\[
  \mu:\mathcal E\longrightarrow\mathcal E\otimes\mathcal C
\]
satisfying axioms \textup{(CM 1)} and \textup{(CM 2)}
of \hyperref[par:sga3-VIB-11.6]{11.6}.

We say that \(\mathcal C\) \textup{(respectively,
\(\mathcal E\))} is a \emph{quasi-coherent coalgebra}
\textup{(respectively, a \emph{quasi-coherent comodule})} if
it is a quasi-coherent \(\mathcal O_S\)-module.

Let \(A\) be a commutative ring and let \(C\) be an
\(A\)-coalgebra. Then
\[
  C^\vee=\operatorname{Hom}_A(C,A)
\]
is an \(A\)-algebra. We denote by
\[
  \operatorname{ev}:C\otimes_A C^\vee\longrightarrow A
\]
the natural evaluation map.

\medskip
\phantomsection\label{lem:sga3-VIB-11.8}%
\noindent\textbf{Lemma 11.8.---}
\emph{Let \(C\) be an \(A\)-coalgebra, \(V\) a
\(C\)-comodule, and \(M\) an \(A\)-submodule of \(V\).
Suppose that \(C\) is a projective \(A\)-module.}%
\footnote{\label{note:sga3-VIB-113}Editor's note: in the
original, \(C\) is assumed to be a free \(A\)-module. The
generalization to the case where \(C\) is a projective
\(A\)-module, pointed out by J.-P.\ Serre, is mentioned in
\hyperref[rem:sga3-VIB-11.10.1]{Remark 11.10.1}. We have incorporated this generalization
here and in \hyperref[prop:sga3-VIB-11.9]{11.9}, and have expanded the proofs
accordingly.}
\emph{Let \(c(M)\) be the image of the morphism of
\(A\)-modules}
\phantomsection\label{form:sga3-VIB-11.8-theta}%
\[
  \theta:
  M\otimes_A C^\vee
  \xrightarrow{\ \mu\otimes\mathrm{id}\ }
  V\otimes_A C\otimes_A C^\vee
  \xrightarrow{\ \mathrm{id}\otimes\operatorname{ev}\ }
  V.
\]
\emph{Then \(c(M)\) is the smallest subcomodule of \(V\)
containing \(M\), and it is a finitely generated
\(A\)-module if \(M\) is. We call \(c(M)\) the
\emph{subcomodule generated by \(M\)}.}

\emph{Moreover, for every homomorphism of rings
\(A\to A'\), if \(M'\) denotes the image of
\(M\otimes_AA'\) in \(V'=V\otimes_AA'\), then \(c(M')\) is
the image of \(c(M)\otimes_AA'\) in \(V'\). Thus the
formation of \(c(M)\) commutes with base change.}

First, \(M\subset c(M)\) by \textup{(CM 2)}. If \(N\) is a
subcomodule of \(V\) containing \(M\), then
\(\mu(M)\subset N\otimes C\), and hence \(c(M)\subset N\).

By hypothesis, \(C\) is a direct summand of a free
\(A\)-module \(L\) with basis \((e_i)_{i\in I}\). Let
\(\varphi_i\) denote the restriction to \(C\) of the linear
form \(e_i^*\), defined by
\[
  e_i^*(e_j)=\delta_{ij}.
\]
Let \(x\in M\). We may write
\phantomsection\label{form:sga3-VIB-11.8-eq1}%
\[
  \tag{1}
  \mu(x)=\sum_{i\in J}x_i\otimes e_i,
\]
where \(x_i\in V\) and \(J\) is a finite subset of \(I\).
Then
\[
  x_i=\theta(x\otimes\varphi_i)\in c(Ax),
\]
and consequently
\[
  c(Ax)=\sum_{i\in J}Ax_i.
\]
Since \(C\) is a direct summand of \(L\), say
\(L=C\oplus R\), whence
\[
  V\otimes L=(V\otimes C)\oplus(V\otimes R),
\]
we obtain
\phantomsection\label{form:sga3-VIB-11.8-intersection}%
\[
  (c(Ax)\otimes L)\cap(V\otimes C)=c(Ax)\otimes C.
\]
Consequently, \(\mu(x)\) may also be written in the form
\phantomsection\label{form:sga3-VIB-11.8-eq2}%
\[
  \tag{2}
  \mu(x)=\sum_{j\in J}x_j\otimes b_j,
\]
with \(b_j\in C\). Write
\[
  \Delta(b_j)=\sum_{i\in I}b_{ij}\otimes e_i,
  \qquad b_{ij}\in C.
\]
Applying \(\mu\otimes\mathrm{id}\) to
\hyperref[form:sga3-VIB-11.8-eq1]{\textup{(1)}} and
\(\mathrm{id}\otimes\Delta\) to
\hyperref[form:sga3-VIB-11.8-eq2]{\textup{(2)}}, and using
axiom \textup{(CM 1)}, we obtain, for every \(i\in J\),
\phantomsection\label{form:sga3-VIB-11.8-coaction-generators}%
\[
  \mu(x_i)=\sum_{j\in J}x_j\otimes b_{ij}
    \in c(Ax)\otimes C.
\]
This proves that \(c(M)\) is a subcomodule of \(V\), and
therefore that it is the smallest subcomodule of \(V\)
containing \(M\).

It is clear that \(c(M)\) is a finitely generated
\(A\)-module whenever \(M\) is: if
\[
  M=Ax_1+\cdots+Ax_n,
  \qquad
  \mu(x_k)=\sum_i x_{ik}\otimes e_i,
\]
then the \(x_{ik}\), for \(k=1,\ldots,n\) and \(i\) ranging
over a finite subset of \(I\), generate \(c(M)\).

Finally, let \(A\to A'\) be a homomorphism of rings, and let
\(M'\) be the image of \(M\otimes_AA'\) in
\(V'=V\otimes_AA'\). Then \(c(M')\), respectively the image
of \(c(M)\otimes_AA'\) in \(V'\), is the image of the
morphism \(\theta'\) below, respectively of the composite
\(\theta'\circ\tau\):
\phantomsection\label{form:sga3-VIB-11.8-base-change-theta}%
\[
  M\otimes_AA'\otimes_A C^\vee
  \xrightarrow{\ \tau\ }
  M\otimes_A\operatorname{Hom}_A(C,A')
  \xrightarrow{\ \theta'\ }V'.
\]
These two morphisms have the same image. Indeed, let
\(\psi\in\operatorname{Hom}_A(C,A')\) and \(x\in M\), and
write
\[
  \mu(x)=\sum_{i\in J}x_i\otimes e_i.
\]
Then
\[
  \theta'(x\otimes\psi)
    =\sum_{i\in J}\psi(p(e_i))x_i
\]
is the image under \(\theta'\circ\tau\) of the element
\[
  \sum_{i\in J}x\otimes\psi(p(e_i))\otimes\varphi_i
    \in M\otimes_AA'\otimes_A C^\vee.
\]
This proves the lemma.

We also have the following proposition.

\medskip
\phantomsection\label{prop:sga3-VIB-11.8-bis}%
\noindent\textbf{Proposition 11.8.bis.---}%
\footnote{\label{note:sga3-VIB-114}Editor's note: this
proposition, taken from \textup{[Se68]}, Section 1.5,
Proposition 2, has been added.}
\emph{Let \(A\) be a Noetherian ring, let \(C\) be an
\(A\)-coalgebra flat over \(A\), let \(V\) be a
\(C\)-comodule, and let \(M\) be a finitely generated
\(A\)-submodule of \(V\). Then there is a subcomodule \(W\)
of \(V\), finitely generated over \(A\), that contains
\(M\).}

Indeed, since \(M\) is finitely generated, so is
\(\Delta_V(M)\). Hence there is a finitely generated
\(A\)-submodule \(M'\) of \(V\) such that
\[
  \Delta_V(M)\subset M'\otimes_A C.
\]
Let \(\pi:V\to V/M'\) be the projection, put
\[
  \overline{\Delta}_V
    =(\pi\otimes\mathrm{id}_C)\Delta_V,
\]
and set
\phantomsection\label{form:sga3-VIB-11.8-bis-W}%
\[
  W=\{x\in V\mid\Delta_V(x)\in M'\otimes_A C\}
    =\ker\overline{\Delta}_V.
\]
This is an \(A\)-submodule of \(V\) that contains \(M\) and
is contained in \(M'\), because
\[
  x=(\mathrm{id}_V\otimes\varepsilon)\Delta_V(x).
\]
It is therefore finitely generated over \(A\). Moreover,
\[
  (\overline{\Delta}_V\otimes\mathrm{id}_C)\Delta_V
    =(\pi\otimes\Delta_C)\Delta_V
\]
vanishes on \(W\); in other words, \(\Delta_V(W)\) is
contained in the kernel \(K\) of
\(\overline{\Delta}_V\otimes\mathrm{id}_C\). Since \(C\) is
flat over \(A\), one has
\[
  K=W\otimes_A C.
\]
Thus \(W\) is a subcomodule of \(V\).

\medskip
\phantomsection\label{lem:sga3-VIB-11.8.1}%
\noindent\textbf{Lemma 11.8.1.---}%
\footnote{\label{note:sga3-VIB-115}Editor's note: this lemma,
which is used in the proof of \SGARefLink{prop:sga3-VIB-11.9}{11.9}, has been added.}
\emph{Let \(C\) be an \(A\)-coalgebra, let \(V\) be a
\(C\)-comodule, let \(M\) be an \(A\)-submodule of \(V\),
and let \(f:A\to A'\) be a faithfully flat homomorphism of
rings. Suppose that \(C'=C\otimes_AA'\) is a projective
\(A'\)-module.}
\begin{enumerate}
\item[\textup{(i)}]\phantomsection\label{concl:sga3-VIB-11.8.1-i}
  \emph{There is a smallest subcomodule \(t(M)\) of \(V\)
  containing \(M\), and \(t(M)\) is a finitely generated
  \(A\)-module if \(M\) is. Moreover, the formation of
  \(t(M)\) commutes with base change.}

\item[\textup{(ii)}]\phantomsection\label{concl:sga3-VIB-11.8.1-ii}
  \emph{More precisely, \(C\) is a projective
  \(A\)-module, and \(t(M)=c(M)\).}
\end{enumerate}

\emph{Proof.}
For \hyperref[concl:sga3-VIB-11.8.1-ii]{\textup{(ii)}},
\textup{[RG71]} \textup{(see \hyperref[prop:sga3-VIB-11.8.2]{Proposition 11.8.2} below)}
shows that \(C\) is a projective \(A\)-module. We may
therefore apply \hyperref[lem:sga3-VIB-11.8]{Lemma 11.8}: \(c(M)\) is the smallest
subcomodule of \(V\) containing \(M\), it is a finitely
generated \(A\)-module if \(M\) is, and its formation
commutes with base change.

To avoid an anachronism \textup{(\([RG71]\) postdates
SGA 3)}, let us sketch a direct proof of
\hyperref[concl:sga3-VIB-11.8.1-i]{\textup{(i)}}. Since
\(A\to A'\) is flat,
\[
  M'=M\otimes_AA'
\]
is an \(A'\)-submodule of \(V'=V\otimes_AA'\). Since
\(C'\) is a projective \(A'\)-module, \(c(M')\) is the
smallest subcomodule of \(V'\) containing \(M'\). Denote by
\(V'\otimes_AA'\) and \(A'\otimes_AV'\) the two
\(A'\otimes_AA'\)-comodule structures on
\[
  V''=V'\otimes_{A'}(A'\otimes_AA')
\]
obtained by the two base changes
\[
  A'\rightrightarrows A'\otimes_AA',
  \qquad
  a'\longmapsto a'\otimes1,
  \quad
  a'\longmapsto1\otimes a'.
\]
The \(A'\)-comodule \(V'\) carries an isomorphism of
\(A'\otimes_AA'\)-comodules
\phantomsection\label{form:sga3-VIB-11.8.1-descent-phi}%
\[
  \varphi:
  V'\otimes_AA'\xrightarrow{\ \sim\ }A'\otimes_AV',
  \qquad
  (x\otimes a')\otimes b'
    \longmapsto b'\otimes(x\otimes a'),
\]
which is a descent datum, that is,
\[
  \varphi_{31}=\varphi_{32}\circ\varphi_{21}.
\]

Since \(M'=M\otimes_AA'\), the isomorphism \(\varphi\)
takes \(M'\otimes_AA'\) onto \(A'\otimes_AM'\), and hence
takes \(c(M'\otimes_AA')\) onto \(c(A'\otimes_AM')\).
Because the formation of \(c(M')\) commutes with base
change,
\[
  c(M'\otimes_AA')=c(M')\otimes_AA',
  \qquad
  c(A'\otimes_AM')=A'\otimes_Ac(M').
\]
Thus
\phantomsection\label{form:sga3-VIB-11.8.1-descent-subcomodule}%
\[
  \varphi\bigl(c(M')\otimes_AA'\bigr)
    =A'\otimes_Ac(M'),
\]
and \(\varphi\) equips \(c(M')\) with a descent datum. By
fpqc descent, there is a unique subcomodule \(t(M)\) of
\(V\) such that
\[
  c(M')=t(M)\otimes_AA'.
\]
It contains \(M\), since \(t(M)\otimes_AA'\) contains
\(M'\). Moreover, if \(N\) is a subcomodule of \(V\)
containing \(M\), then \(N\) contains \(t(M)\), because
\(N\otimes_AA'\) contains
\[
  c(M')=t(M)\otimes_AA'.
\]
Hence \(t(M)\) is the smallest subcomodule of \(V\)
containing \(M\).

Finally, let \(A\to B\) be a homomorphism of rings. Put
\[
  B'=B\otimes_AA',
\]
and let \(M_B\), respectively \(M'_{B'}\), be the image of
\(M\otimes_AB\) in \(V_B=V\otimes_AB\), respectively of
\(M'\otimes_{A'}B'\) in
\(V'\otimes_{A'}B'=V\otimes_AB'\). Then
\[
  M_B\otimes_BB'=M'_{B'}.
\]
On the one hand, the preceding construction, applied to
\(C_B\) and \(B\to B'\), gives
\phantomsection\label{form:sga3-VIB-11.8.1-base-change-first}%
\[
  c(M_B\otimes_BB')
    =t(M_B)\otimes_BB'
    =t(M_B)\otimes_AA'.
\]
On the other hand, because the formation of \(c(M')\)
commutes with base change, \(c(M'_{B'})\) is the image in
\[
  V'\otimes_{A'}B'=V\otimes_AB\otimes_AA'
\]
of
\phantomsection\label{form:sga3-VIB-11.8.1-base-change-second}%
\[
  c(M')\otimes_{A'}B'
    =t(M)\otimes_AB\otimes_AA'.
\]
It follows that \(t(M_B)\) is the image of
\(t(M)\otimes_AB\) in \(V_B\).

\medskip
\phantomsection\label{prop:sga3-VIB-11.8.2}%
\noindent\textbf{Proposition 11.8.2 \textup{(Gruson--Raynaud)}.---}
\emph{Let \(f:A\to A'\) be a faithfully flat homomorphism.
Then \(f\) ``descends projectivity'': if \(M\) is an
\(A\)-module and \(M\otimes_AA'\) is a projective
\(A'\)-module, then \(M\) is a projective \(A\)-module.}

Indeed, by \textup{[RG71]}, II, \SGARefLink{lem:sga3-VIA-2.5.1}{2.5.1}, \(f\) descends the
Mittag--Leffler condition, and therefore, by
\emph{loc.\ cit.}, II, 3.1.3, \(f\) descends
projectivity.%
\footnote{\label{note:sga3-VIB-116}Editor's note: let us
also point out that assertion II, \SGARefLink{prop:sga3-VIA-2.5.2}{2.5.2} of \emph{loc.\ cit.},
which is more general than II, \SGARefLink{lem:sga3-VIA-2.5.1}{2.5.1}, is corrected in
\textup{[Gr73]}; this does not affect the case of faithfully
flat homomorphisms.}

\medskip
\phantomsection\label{prop:sga3-VIB-11.9}%
\noindent\textbf{Proposition 11.9.---}%
\footnote{\label{note:sga3-VIB-117}Editor's note: as
indicated in editor's \hyperref[note:sga3-VIB-112]{note 112}, we have rewritten the
statement for an \(\mathcal O_S\)-coalgebra \(\mathcal C\)
\textup{(rather than for an \(S\)-group \(G\) satisfying the
hypotheses in \hyperref[cor:sga3-VIB-11.10]{Corollary 11.10})}. We have also expanded the
proof \textup{(the original merely said: ``\(\cdots\) the
proposition follows from \hyperref[lem:sga3-VIB-11.8]{Lemma 11.8}'')}.}
\emph{Let \(\mathcal C\) be an \(\mathcal O_S\)-coalgebra,
let \(\mathcal E\) be a \(\mathcal C\)-comodule, and let
\(\mathcal F\) be an \(\mathcal O_S\)-submodule of
\(\mathcal E\), all of them quasi-coherent. Suppose that
\(S\) is covered by affine open subsets
\(U_\alpha=\operatorname{Spec}A_\alpha\), and that for each
\(\alpha\) there is a faithfully flat homomorphism of rings
\(A_\alpha\to A'_\alpha\) such that}
\[
  \Gamma(U_\alpha,\mathcal C)\otimes_{A_\alpha}A'_\alpha
\]
\emph{is a projective \(A'_\alpha\)-module.}%
\textsuperscript{\textup{(*)}}
\begingroup
\renewcommand{\thefootnote}{*}%
\footnotetext{This holds, for example, when
\(\mathcal C=\mathcal A(G)\), where \(G\) is a reductive
\(S\)-group, as we shall see in \SGARefLink{sga3:xxii:cor:5.7.8}{Exposé XXII, 5.7.8}.}
\endgroup

\emph{Then there is a smallest quasi-coherent subcomodule
\(t(\mathcal F)\) of \(\mathcal E\) containing
\(\mathcal F\), and \(t(\mathcal F)\) is a finitely
generated \(\mathcal O_S\)-module if \(\mathcal F\) is.
Moreover, for every base change \(S'\to S\), if
\(\mathcal F'\) denotes the image of
\(\mathcal F\otimes_{\mathcal O_S}\mathcal O_{S'}\) in}
\[
  \mathcal E'
    =\mathcal E\otimes_{\mathcal O_S}\mathcal O_{S'},
\]
\emph{then \(t(\mathcal F')\) is the image of
\(t(\mathcal F)\otimes_{\mathcal O_S}\mathcal O_{S'}\) in
\(\mathcal E'\). In other words, the formation of
\(t(\mathcal F)\) commutes with base change.}

\emph{Proof.}
For each \(\alpha\), let \(T_\alpha\) be the
\(A_\alpha\)-module
\[
  T_\alpha=t\bigl(\Gamma(U_\alpha,\mathcal F)\bigr).
\]
By \hyperref[concl:sga3-VIB-11.8.1-i]{11.8.1}\textup{\SGARefLink{concl:sga3-VIB-11.8.1-i}{(i)}}, the associated
\(\mathcal O_{U_\alpha}\)-module
\(\widetilde T_\alpha\) is the smallest quasi-coherent
subcomodule of \(\mathcal E|_{U_\alpha}\) containing
\(\mathcal F|_{U_\alpha}\); and it is a finitely generated
\(\mathcal O_{U_\alpha}\)-module if
\(\mathcal F|_{U_\alpha}\) is.

For every \(\alpha,\beta\), put
\[
  U_{\alpha\beta}=U_\alpha\cap U_\beta.
\]
Since the construction of \(t(M)\) commutes with base
change, there are canonical isomorphisms of
\(\mathcal O_{U_{\alpha\beta}}\)-modules
\phantomsection\label{form:sga3-VIB-11.9-gluing}%
\[
  \varphi_{\alpha\beta}:
  \widetilde T_\beta
    \otimes_{A_\beta}\mathcal O_{U_{\alpha\beta}}
  \xrightarrow{\ \sim\ }
  \widetilde T_\alpha
    \otimes_{A_\alpha}\mathcal O_{U_{\alpha\beta}}.
\]
They satisfy the cocycle condition
\[
  \varphi'_{\alpha\gamma}
    =\varphi'_{\alpha\beta}\circ\varphi'_{\beta\gamma},
\]
where \(\varphi'_{\alpha\gamma}\), respectively the other
maps, denotes the restriction of
\(\varphi_{\alpha\gamma}\), respectively the other maps, to
\(U_\alpha\cap U_\beta\cap U_\gamma\).

Consequently, the \(\widetilde T_\alpha\) glue to a
quasi-coherent subcomodule \(t(\mathcal F)\) of
\(\mathcal E\) containing \(\mathcal F\). We leave it to the
reader to verify that \(t(\mathcal F)\) is the smallest
quasi-coherent subcomodule of \(\mathcal E\) containing
\(\mathcal F\), and that its formation commutes with base
change.

\medskip
\phantomsection\label{def:sga3-VIB-11.9.1}%
\noindent\textbf{Definition 11.9.1.---}%
\footnote{\label{note:sga3-VIB-118}Editor's note: this
definition, taken from the bottom of page 82 of
\textup{[RG71]}, has been added. Thus the hypothesis in
\hyperref[prop:sga3-VIB-11.9]{Proposition 11.9} is that the coalgebra \(\mathcal C\) be a
locally projective \(\mathcal O_S\)-module, and we have used
this terminology in \hyperref[cor:sga3-VIB-11.10]{Corollary 11.10}.}
\emph{Let \(S\) be a scheme and let \(\mathcal P\) be a
quasi-coherent \(\mathcal O_S\)-module. The following
conditions are equivalent:}
\begin{enumerate}
\item[\textup{(i)}]\phantomsection\label{concl:sga3-VIB-11.9.1-i}
  \emph{For every affine open subset \(U\) of \(S\), the
  module \(\Gamma(U,\mathcal P)\) is a projective
  \(\mathcal O_S(U)\)-module.}

\item[\textup{(ii)}]\phantomsection\label{concl:sga3-VIB-11.9.1-ii}
  \emph{There is a covering \((U_\alpha)\) of \(S\) by
  affine open subsets such that every
  \(\Gamma(U_\alpha,\mathcal P)\) is a projective
  \(\mathcal O_S(U_\alpha)\)-module.}

\item[\textup{\SGAInlineRefTarget{def:sga3-VIB-11.9.1:item:iii}(iii)}]\phantomsection\label{concl:sga3-VIB-11.9.1-iii}
  \emph{There is a covering \((U_\alpha)\) of \(S\) by
  affine open subsets, and faithfully flat homomorphisms of
  rings}
  \[
    A_\alpha=\mathcal O_S(U_\alpha)\longrightarrow A'_\alpha,
  \]
  \emph{such that, for every \(\alpha\),}
  \[
    \Gamma(U_\alpha,\mathcal P)
      \otimes_{A_\alpha}A'_\alpha
  \]
  \emph{is a projective \(A'_\alpha\)-module.}
\end{enumerate}

It is clear that
\[
  \textup{(i)}\Longrightarrow\textup{(ii)}
  \Longrightarrow\textup{(iii)}.
\]
Conversely, if \textup{\SGARefLink{def:sga3-VIB-11.9.1:item:iii}{(iii)}} holds, \hyperref[prop:sga3-VIB-11.8.2]{11.8.2} implies that
every \(\Gamma(U_\alpha,\mathcal P)\) is a projective
\(\mathcal O_S(U_\alpha)\)-module, and hence
\textup{(ii)} holds.

Finally, suppose that \textup{\SGARefLink{concl:sga3-VIB-11.9.1-ii}{(ii)}} holds, and let
\[
  V=\operatorname{Spec}A
\]
be an arbitrary affine open subset. It is covered by
finitely many affine open subsets
\[
  V_1,\ldots,V_n,
  \qquad
  V_i=\operatorname{Spec}A_i,
\]
each contained in at least one \(V\cap U_\alpha\), so that
\(\Gamma(V_i,\mathcal P)\) is a projective
\(A_i\)-module. Put
\[
  A'=A_1\times\cdots\times A_n.
\]
Then \(A\to A'\) is faithfully flat and
\[
  \Gamma(V,\mathcal P)\otimes_AA'
\]
is a projective \(A'\)-module. Therefore \SGARefLink{prop:sga3-VIB-11.8.2}{11.8.2} implies that
\(\Gamma(V,\mathcal P)\) is a projective \(A\)-module.

When these equivalent conditions hold, we say that
\(\mathcal P\) is a \emph{locally projective}
\(\mathcal O_S\)-module.

\medskip
\phantomsection\label{cor:sga3-VIB-11.10}%
\noindent\textbf{Corollary 11.10.---}
\emph{Let \(S\) be a quasi-compact and quasi-separated
scheme, let \(G\) be an \(S\)-group, and let \(\rho\) be a
linear action of \(G\) on a quasi-coherent
\(\mathcal O_S\)-module \(\mathcal E\). Suppose that:}
\begin{enumerate}
\item[\textup{(i)}]\phantomsection\label{hyp:sga3-VIB-11.10-i}
  \emph{\(G\) satisfies one of the following conditions:}
  \begin{enumerate}
  \item[\textup{(a)}]\phantomsection\label{hyp:sga3-VIB-11.10-i-a}
    \emph{\(G\) is affine and flat over \(S\);}

  \item[\textup{(b)}]\phantomsection\label{hyp:sga3-VIB-11.10-i-b}
    \emph{\(G\) is flat, quasi-compact, and quasi-separated
    over \(S\), and \(\mathcal E\) is flat;}
  \end{enumerate}

\item[\textup{(ii)}]\phantomsection\label{hyp:sga3-VIB-11.10-ii}
  \emph{\(\mathcal A(G)\) is a locally projective
  \(\mathcal O_S\)-module.}
\end{enumerate}
\emph{Then \(\mathcal E\) is the direct limit of an
increasing filtered family of finitely generated
quasi-coherent \(\mathcal O_S\)-submodules of
\(\mathcal E\) that are stable under \(G\).}

By hypothesis
\hyperref[hyp:sga3-VIB-11.10-i]{\textup{(i)}} and \hyperref[prop:sga3-VIB-11.6.1]{11.6.1},
\(\mathcal E\) is endowed with an
\(\mathcal A(G)\)-comodule structure. On the other hand,
since \(S\) is quasi-compact and quasi-separated,
\(\mathcal E\) is the direct limit of its finitely
generated quasi-coherent submodules
\textup{(EGA I, 9.4.9 and EGA IV\(_1\), 1.7.7)}.
Consequently, the corollary follows from \hyperref[prop:sga3-VIB-11.9]{Proposition 11.9},
applied to the coalgebra \(\mathcal A(G)\).

We also have the following proposition.
